% ===================================================================
%  పట్టిక సంఖ్య 8 — కోత.
%
%  Traduction, faite en 2026, en telougou d'aujourd'hui, sur l'ido de
%  texto/io. Regles de la colonne : voir 10-pattika-01.tex. Deux
%  scenes, deux numerotations : les cles portent c1 et c2. Le
%  fac-simile compose « (13) » sans sa parenthese ouvrante ; la
%  traduction ecrit la reference bien formee, comme le font deja les
%  colonnes hindi et marathe.
%
%  LE FLEAU MANQUE UNE TROISIEME FOIS, ET LES TROIS LANGUES ONT
%  ECRIT LA MEME CHOSE. L'egyptien n'avait pas de mot parce que
%  l'Egypte bat au نورج : il ecrivait عصي الدراس, « les batons de
%  battage ». Le marathi n'en avait pas parce que le Maharashtra bat
%  en faisant pietiner les bœufs : मळणीच्या काठ्या. Le telougou
%  n'en a pas pour la meme raison que le marathi, et sur un terrain
%  qui n'est meme pas de la meme famille de langues : నూర్పిడి
%  కర్రలు. Trois colonnes, trois familles — semitique, indo-aryenne,
%  dravidienne —, un seul trou, et la meme periphrase mot pour mot.
%  Le trou est dans la CHOSE, et il se voit d'autant mieux que les
%  langues n'ont rien en commun.
%
%  ET LE RESTE DE LA MOISSON EST ENTIEREMENT TELOUGOU, ce qui rend
%  le trou plus net encore :
%
%    కళ్ళం        l'aire          పన (31)      la gerbe
%    సాన (30)     la pierre a aiguiser
%    కొడవలి (12)  la faucille     పెద్ద కొడవలి (25) la faux
%    కంకులు (29)  les epis        గడ్డివాము (44) la meule
%    బీడు (9)     le paturage     ఊట (20)      la source
%    వాగు (16)    le ruisseau     అడవి (49)    le bois
%    గానుగ (28)   le pressoir     గెల          la grappe
%    తీగ (39)     le cep          పీపా (20)    le tonneau
%    గాలం (44)    l'hamecon       ఎర           l'appat
%    వల (36)      le filet        పంజరం (34)   la cage
%    సొరకాయ బుర్ర (14) la gourde — la calebasse sechee ou l'on porte
%                 l'eau, qui est exactement l'objet du livret
%
%  TROIS OISEAUX QUE JE NE SAIS PAS NOMMER, et je prefere le dire.
%  L'hirondelle (deja au tableau 3), la bergeronnette et l'alouette
%  n'ont pas de forme telougoue dont je sois sur. Le telougou a des
%  noms pour la plupart des oiseaux de ce livret — పిచ్చుక le
%  moineau, కాకి la corneille, నెమలి le paon, కోడిపుంజు le coq —,
%  mais pour ces trois-la j'ecris l'emprunt ou la description, et le
%  trou est note ici plutot que masque. Inventer une forme aurait ete
%  pire : le lecteur aurait cru a un mot.
%
%  ET DEUX EMPRUNTS DE PLANTES, MEME REGLE QU'AU MARRON CHAUD : le
%  bleuet et la digitale ne poussent pas en Andhra. Le coquelicot,
%  lui, tombe juste — గసగసాలు, le pavot, pousse en Inde depuis
%  toujours et donne son nom au meme grain.
%
%  ECRIT AVEC LA LISTE DE parigi.py, ET DEUX INVERSIONS ONT PASSE
%  QUAND MEME — les deux absentes de la liste, et c'est instructif :
%
%   * « la planajo (39) UBE flugeskas la alaudi (41) » : « ube » est
%     l'adverbe relatif de lieu, de la meme famille que « qua », et
%     le motif de l'outil l'avait oublie. Vrai trou, comble.
%   * « la fisho (47) snapas l'angelo quik KANDE ili plunjas sua filo
%     (46) » : « kande » introduit une circonstancielle, pas un
%     modificateur de nom. Ce n'est PAS un rapport modificateur-tete,
%     et l'outil a raison de ne pas le signaler — l'inversion venait
%     de moi, qui avais mis la ligne avant le poisson. On n'ajoute
%     donc ni « kande » ni « dum ».
%
%  Le premier cas elargit l'outil ; le second lui donne raison. Il
%  fallait les distinguer plutot que d'ajouter les deux.
% ===================================================================

%%K t08-noto-1 noto
\VUnotes{143.2mm}{%
\textit{(*) ఎందుకంటే ఈ మాట కచ్చితమైనది కాదు : ఇది అవిసె కోతా, ద్రాక్ష కోతా, లేక ఆపిల్ కోతా?}
\textit{Mondo XIV, పుట 242 లో సూచించిన «} \VUgras{mesono}
\textit{» అనే మాట వాడటం మేలేమో. కానీ ఇప్పటివరకు ఆ కొత్త ధాతువును అకాడమీ స్వీకరించలేదు, ఇడో ఆనవాయితీ నియమాల ప్రకారం అది చర్చ కోసం వేచి ఉంది.}%
}

%%K t08-apar-1 apar
\VUsaut{5.0mm}
\VUtitre{15.0pt}{60}{పట్టిక సంఖ్య 8}
\VUsaut{3.0mm}
\VUcentre{13.2pt}{90}{\textit{మొదటి రంగం.}}
\VUsaut{3.0mm}
\VUpk{10.2pt}{40}{కోత (*).}
\VUsaut{4.0mm}
\VUpk{10.2pt}{40}{పొలం రూపురేఖలు.}
\VUsaut{3.0mm}

%%K t08-c1-01-1 p
1. --- \VUgras{వేసవితో} పాటు పొలం రూపురేఖలు మరోసారి మారాయి. చాలుకు అప్పగించిన విత్తనం
మొలకెత్తింది; నేలలోంచి పచ్చని మొక్క పైకి వచ్చి కంకి పట్టింది. గింజ కట్టింది, ఆ తర్వాత పండింది,
ఇప్పుడు కోసుకోవడం తప్ప ఏమీ మిగల్లేదు.

%%K t08-c1-02-1 p
2. --- \VUgras{పొలపు పూలు} విచ్చుకున్నాయి. \VUgras{కార్న్‌ఫ్లవర్} \textsuperscript{(1)},
\VUgras{గసగసాల పూలు} \textsuperscript{(2)} మరియు \VUgras{ఫాక్స్‌గ్లవ్} \textsuperscript{(3)}
గోధుమ చేల ఒకేలాంటి రంగులో తమ తాజా \VUgras{రేకుల} ఆనందకరమైన వైరుధ్యాన్ని కలుపుతున్నాయి.

%%K t08-c1-03-1 p
3. --- పండ్ల చెట్లకు ఆకులు వచ్చాయి; పండు పండింది. చిన్న \VUgras{చెర్రీ చెట్టుకు}
\textsuperscript{(4)} అందమైన \VUgras{చెర్రీలు} \textsuperscript{(5)} బాగా కాశాయి, వాటిని
పిల్లలు ఎంతో ఆనందంగా కోసి తింటారు.

%%K t08-c1-04-1 p
4. --- ఇక నుంచి పశువులు రోజంతా బయటి గాలిలో గడపవచ్చు.

%%K t08-c1-04-2 p
\VUgras{గొర్రెపిల్లలు} \textsuperscript{(6)} పెద్దవై దట్టమైన ఉన్నితో అందమైన \VUgras{గొర్రెలు}
\textsuperscript{(7)} మరియు అందమైన \VUgras{పొట్టేళ్ళు} \textsuperscript{(8)} అయ్యాయి. డైసీ
పూలతో రంగురంగులైన \VUgras{మేతబీడులోని} \textsuperscript{(9)} \VUgras{గడ్డిని} అవి ప్రశాంతంగా
మేస్తాయి.

%%K t08-c1-04-3 p
త్వరలో \VUgras{ఆ} పశువులను కత్తిరిస్తారు, వాటి \VUgras{ఉన్ని} తీసేందుకు; దానితో మన బట్టల గుడ్డ
తయారవుతుంది.

%%K t08-tit-2 sub
\VUsaut{5.0mm}
\VUpk{10.2pt}{40}{గొర్రెల కాపరి.}
\VUsaut{3.4mm}

%%K t08-c1-05-1 p
5. --- \VUgras{గొర్రెల కాపరి} \textsuperscript{(10)} వాటిని పెద్దగా పట్టించుకోనట్టున్నాడు.
క్షితిజం మీదుగా వెళ్ళే తుఫాను గురించే అతనికి బెంగ; \VUgras{గొర్రెల కావలివాడు}
\textsuperscript{(13)} మరింత నిర్లక్ష్యంగా కనిపిస్తాడు — అతను తన \VUgras{సొరకాయ బుర్రను}
\textsuperscript{(14)} పెదవుల దగ్గరికి తెచ్చుకుంటున్నాడు.

%%K t08-c1-06-1 p
6. --- ఎండ భరించలేనట్టుగా ఉంది; ఎందుకంటే \VUgras{కుక్క} \textsuperscript{(15)} కూడా దాహం
తీర్చుకునేందుకు వస్తోంది; అది \VUgras{వాగులోని} \textsuperscript{(16)} చల్లని నీళ్ళు నాకుతూ,
ఒడ్డు గడ్డిలో దాక్కున్న పాపం \VUgras{కప్పను} \textsuperscript{(17)} భయపెడుతోంది. అది స్వచ్ఛమైన
నీళ్ళలోకి దూకుతుంది, అక్కడ \VUgras{తామరలు} \textsuperscript{(18)} పూశాయి.

%%K t08-c1-07-1 p
7. --- ఒక \VUgras{తోక ఊపే పిట్ట} \textsuperscript{(19)} \VUgras{ఊట} \textsuperscript{(20)}
దగ్గర స్నానం చేస్తోంది — ఆ ఊట రెండు బండల మధ్యనుంచి ఉబుకుతోంది — ఆ తర్వాత
\VUgras{ముళ్ళ పొదల కింద} \textsuperscript{(21)}, \VUgras{హాథార్న్ పొదల కింద}
\textsuperscript{(22)} దాక్కుంటోంది.

%%K t08-tit-3 sub
\VUsaut{5.0mm}
\VUpk{10.2pt}{40}{కోత.}
\VUsaut{3.4mm}

%%K t08-c1-08-1 p
8. --- పల్లె పిల్లలకు గోధుమ కోత చాలా సరదా కాలం. వాళ్ళు ఆనందంగా \VUgras{దొర్లుతారు}
\textsuperscript{(24)}, \VUgras{గడ్డి కుప్పలపై} \textsuperscript{(23)}. వాళ్ళు
\VUgras{కోతకూలీలకు} \textsuperscript{(26)} సాయం చేస్తారు, ఆ కూలీలు \VUgras{గోధుమ చేను}
\textsuperscript{(27)} కోస్తుంటే — తమ \VUgras{పెద్ద కొడవళ్ళతో} \textsuperscript{(25)}, అవి
\VUgras{సానపై} \textsuperscript{(30)} బాగా పదును పెట్టినవి — పిల్లలు గోధుమ పోగుచేసి
\VUgras{పనలుగా} \textsuperscript{(31)} లేదా \VUgras{కట్టలుగా} \textsuperscript{(32)} కడతారు.

%%K t08-c1-09-1 p
9. --- అప్పుడప్పుడూ వాళ్ళలో అందరికంటే పెద్దవాడికి \VUgras{కొడవలితో} \textsuperscript{(12)} ఆ
\VUgras{బంగారు కాడలు} \textsuperscript{(28)} కోసేందుకు అనుమతిస్తారు — ఆ కాడలు గింజలతో నిండిన
బరువైన \VUgras{కంకులను} \textsuperscript{(29)} మోస్తాయి; చిన్నవాళ్ళు \VUgras{పశువులను}
\textsuperscript{(33)} కనిపెడుతూ — అవి \VUgras{మేతబీడులో} \textsuperscript{(34)} మేస్తాయి,
పెద్ద \VUgras{లిండెన్ చెట్టు} \textsuperscript{(35)} నీడలో — తమ \VUgras{వలతో}
\textsuperscript{(36)} అందమైన \VUgras{సీతాకోకచిలుకలు} పడతారు.

%%K t08-c1-09-2 p
కొందరైతే \VUgras{బండిపైకి} \textsuperscript{(37)} ఎక్కి, \VUgras{గడ్డి పారతో}
\textsuperscript{(38)} తమకు అందించే పనలు సర్దుతారు.

%%K t08-c1-10-1 p
10. --- \VUgras{మైదానమంతా} \textsuperscript{(39)} చైతన్యం నిండింది — ఆ మైదానంలో
\VUgras{లార్క్ పిట్టలు} \textsuperscript{(41)} ఎగరడం మొదలుపెడతాయి. అందరూ కోతలో మునిగి ఉన్నారు.
కానీ పేదవాళ్ళు పాత పనిముట్లే వాడుతూ ఉంటారు — \VUgras{పెద్ద కొడవళ్ళు, కొడవళ్ళు} మరియు
\VUgras{నూర్పిడి కర్రలు} \textsuperscript{(40)} —, ధనిక భూస్వాములు మాత్రం తమ గోధుమలను లేదా తమ
\VUgras{రై ధాన్యాన్ని} \textsuperscript{(43)} \VUgras{కోత యంత్రంతో} \textsuperscript{(42)}
కోయిస్తారు. ఆ తర్వాత పనలను కళ్ళానికి మోయకుండా అక్కడే పెద్ద \VUgras{గడ్డివాములు}
\textsuperscript{(44)} వేస్తారు.

%%K t08-c1-11-1 p
11. --- అప్పుడు \VUgras{నూర్పిడి యంత్రం} \textsuperscript{(47)} తీసుకొస్తారు, దాన్ని
\VUgras{ఆవిరి యంత్రం} \textsuperscript{(45)} నడిపిస్తుంది — \VUgras{పట్టీతో}
\textsuperscript{(46)} —, ఆ విధంగా గింజ \VUgras{గడ్డి} నుంచి వేరవుతుంది, విత్తిన చేను
వదలకుండానే.

%%K t08-tit-4 sub
\VUsaut{5.0mm}
\VUpk{10.2pt}{40}{తుఫాను.}
\VUsaut{3.4mm}

%%K t08-c1-12-1 p
12. --- కోత కాలంలో తరచూ \VUgras{పెద్ద తుఫానులు} \textsuperscript{(53)} వస్తాయి. ముందుగా దూరం
నుంచి \VUgras{ఉరుముల} చప్పుడు వినిపిస్తుంది; \VUgras{పడమటి గాలివాన} \textsuperscript{(54)}
మొదలై, ఆయాసపడుతూ \VUgras{అడవిలోని} \textsuperscript{(49)} అతి లావు చెట్లను వంచుతుంది. నల్లని
\VUgras{మేఘాలు} \textsuperscript{(56)} ఆకాశంపై దాడి చేస్తాయి; \VUgras{మెరుపు}
\textsuperscript{(55)} కళ్ళు మిరుమిట్లు గొలిపే వంకరటింకర గీతలతో వాటిని చీలుస్తుంది.
\VUgras{వాన}, అప్పుడప్పుడూ \VUgras{వడగళ్ళు} పడతాయి, కోతకు సిద్ధమైన పంటను నేలకూలుస్తూ.

%%K t08-c1-13-1 p
13. --- మెరుపు తరచూ ఏదో ఒక కట్టడాన్ని తగలబెడుతుంది. అలాగే గత సంవత్సరం \VUgras{గాలిమర}
\textsuperscript{(50)} కప్పు బూడిదైపోయింది, దాని రెండు \VUgras{రెక్కలు} \textsuperscript{(51)}
విరిగాయి.

%%K t08-c1-14-1 p
14. --- కొన్ని గంటల తర్వాత మళ్ళీ ప్రశాంతత తిరిగి వస్తుంది. క్షితిజంపై ప్రకాశవంతమైన
\VUgras{ఇంద్రధనుస్సు} \textsuperscript{(52)} కనిపిస్తుంది, కొన్ని క్షణాలు ఆగిన ప్రకృతి తన
జీవితాన్ని తిరిగి మొదలుపెడుతుంది. \VUgras{గుడిసెల్లో} \textsuperscript{(48)} ఆశ్రయం పొందిన
పల్లెవాసులు మళ్ళీ పనికి దిగుతారు, ఇప్పుడే పడ్డ నష్టాన్ని ధైర్యంగా సరిదిద్దుకునేందుకు
ప్రయత్నిస్తారు.

%%K t08-tit-5 sub
\VUsaut{6.0mm}
\VUcentre{13.2pt}{90}{\textit{రెండవ రంగం.}}
\VUsaut{3.6mm}

%%K t08-tit-6 sub
\VUpk{10.2pt}{40}{వేట. --- ద్రాక్ష కోత.}
\VUsaut{1.4mm}
\VUpk{10.2pt}{40}{చేపల వేట.}
\VUsaut{4.0mm}

%%K t08-c2-01-1 p
1. --- ప్రాంతాన్ని బట్టి \VUgras{ఆగస్టు} చివరిలో లేదా \VUgras{సెప్టెంబరు} మధ్యలో వేట కాలం
మొదలవుతుంది. సూర్యుడి తొలి కిరణాలు \VUgras{తొండలను} \textsuperscript{(2)} వాటి కన్నాల్లోంచి
బయటికి తీశాయో లేదో, \VUgras{వేటగాడు} \textsuperscript{(5)} బయలుదేరుతాడు — తన \VUgras{కుక్కలతో}
\textsuperscript{(4)}.

%%K t08-c2-02-1 p
2. --- మందపాటి గుడ్డతో చేసిన తన దృఢమైన \VUgras{వేట దుస్తులు} \textsuperscript{(9)} అతను
వేసుకున్నాడు, కాళ్ళకు తోలు \VUgras{కాళ్ళ పట్టీలు} కట్టుకున్నాడు — వాటి సాయంతో అతను ఆ తడి
గడ్డిలో నడవగలడు, అక్కడ \VUgras{గోదురు కప్పలు} \textsuperscript{(1)} దాక్కుంటాయి; తన
\VUgras{తుపాకి} \textsuperscript{(6)} మరియు తన \VUgras{వేట సంచి} \textsuperscript{(10)}
తీసుకున్నాడు, ఇదిగో అతను బయలుదేరాడు.

%%K t08-c2-03-1 p
3. --- వేట ఉత్సాహంలో అతను ద్రాక్షతోట మధ్యకు చేరుకుంటాడు, అక్కడ ద్రాక్ష కోత జోరుగా సాగుతోంది.
ద్రాక్షతోట రైతు పిల్లలు, మామూలుగా దారిలో మేకలు కాస్తారు, ఈ రోజు వాటిని ఎక్కడ పడితే అక్కడ
మేయనిస్తారు; తన స్వేచ్ఛను వాడుకుని పారిపోవడంలో వెనకాడని ముసలి \VUgras{మేకపోతును}
\textsuperscript{(3)} కొయ్యకు కట్టేసి, వాళ్ళూ ఈ సరదాలో తమ వాటా తీసుకుంటారు. ఒకడు
\VUgras{కోత బుట్టలోని} \textsuperscript{(12)} ద్రాక్ష గెల తీసుకుని, \VUgras{రసాన్ని}
\textsuperscript{(14)} ఒక \VUgras{గ్లాసులోకి} \textsuperscript{(13)} పిండి కొత్త రసం (పులియని
ద్రాక్షసారా) చేసే సరదా చేస్తాడు. ఇంకొకడు ఇప్పుడే అల్లిన \VUgras{ఆకుల కిరీటంతో}
\textsuperscript{(15)} తల అలంకరించుకుంటాడు. వాళ్ళ దగ్గర సిగ్గులేని \VUgras{పిచ్చుకలు}
\textsuperscript{(16)} ద్రాక్షలు దొంగిలించేందుకు గానుగ ముందుదాకా వస్తాయి.

%%K t08-c2-04-1 p
4. --- పిల్లలు సరదా పడుతుంటే మగవాళ్ళు పనిచేస్తారు. \VUgras{ద్రాక్ష కోసేవాళ్ళు}
\textsuperscript{(18)} ద్రాక్షలను \VUgras{వీపు బుట్టలతో} \textsuperscript{(17)} నేలమాళిగలోకి
మోసుకెళ్తారు; \VUgras{పీపాలు చేసేవాడు} \textsuperscript{(19)} \VUgras{పీపాలు}
\textsuperscript{(20)} బాగుచేస్తాడు, \VUgras{పలకలు} \textsuperscript{(22)} మరియు
\VUgras{అడుగులు} \textsuperscript{(21)} బాగున్నాయో లేదో చూస్తాడు, విరిగిన
\VUgras{ఇనుప కడియాలు} \textsuperscript{(23)} మారుస్తాడు. \VUgras{పెద్ద తొట్లు}
\textsuperscript{(25)} కూడా అతనే చేశాడు — వాటిలో \VUgras{ద్రాక్షలు} \textsuperscript{(26)}
పులియబెట్టేందుకు పోస్తారు.

%%K t08-c2-05-1 p
5. --- పులియడం పూర్తయ్యాక ద్రాక్షసారా తీసేసి, మిగిలిన పిప్పిని \VUgras{గానుగపై}
\textsuperscript{(28)} పెడతారు; \VUgras{పెద్ద మరతో} \textsuperscript{(29)} కొద్దికొద్దిగా
అదుముతూ రసం పిండుతారు, అది \VUgras{తొట్టిలోకి} \textsuperscript{(32)} కారుతుంది, దానినుంచీ
\VUgras{వైన్} \textsuperscript{(31)} వస్తుంది.

%%K t08-tit-8 sub
\VUsaut{6.0mm}
\VUpk{10.2pt}{40}{బీరు మరియు సైడర్.}
\VUsaut{3.4mm}

%%K t08-c2-06-1 p
6. --- \VUgras{నేలమాళిగ} \textsuperscript{(27)} గోడపై \VUgras{హాప్స్} \textsuperscript{(30)}
కొమ్మ ఒకటి ఎక్కుతోంది. ఇక్కడ అది కేవలం అలంకారపు తీగ, కానీ ద్రాక్షతీగ చాలా అరుదైన కొన్ని
దేశాల్లో \VUgras{బీరు} చేసేందుకు హాప్స్ సాగుచేస్తారు.

%%K t08-c2-07-1 p
7. --- చిన్న \VUgras{ఆపిల్ చెట్టుకు} \textsuperscript{(33)} ఆపిల్స్ బాగా కాశాయి, అవి పండాక
చక్కని \VUgras{సైడర్} ఇస్తాయి. తమ \VUgras{పంజరంలో} \textsuperscript{(34)} \VUgras{కేనరీ}
\textsuperscript{(35)} మరియు \VUgras{గోల్డ్‌ఫించ్} \textsuperscript{(36)} స్వేచ్ఛగా గంతులేసే
పిచ్చుకలవైపు అసూయగా చూస్తాయి.

%%K t08-c2-08-1 p
8. --- కంటిచూపు అందేంతవరకూ, కొండల వాలుపై \VUgras{ద్రాక్ష కొయ్యలు} \textsuperscript{(40)}
వరుసగా నిలబడి ఉన్నాయి, అవి అందమైన \VUgras{ద్రాక్ష తీగలను} \textsuperscript{(39)} మోస్తాయి — ఆ
తీగలకు బరువైన \VUgras{గెలలు} లగడ్డాయి.

%%K t08-c2-08-2 p
ఏడాదంతా \VUgras{ద్రాక్షతోట రైతు} \textsuperscript{(42)} కష్టపడ్డాడు, తన \VUgras{ద్రాక్షతోట}
\textsuperscript{(38)} బాగోగులు చూసేందుకు; ఇప్పుడు చక్కని పంటతో అతని శ్రమకు ఫలం దక్కింది.
\VUgras{ద్రాక్ష కోసే స్త్రీలు} \textsuperscript{(41)} \VUgras{కంచరగాడిదలు}
\textsuperscript{(43)} లాగే బండిపై పెట్టిన తొట్లు నింపుతారు, అందరూ ఆనందంగా ఉన్నారు.

%%K t08-tit-9 sub
\VUsaut{5.0mm}
\VUpk{10.2pt}{40}{చేపల వేట.}
\VUsaut{3.4mm}

%%K t08-c2-09-1 p
9. --- \VUgras{గాలం వేసేవాళ్ళ} \textsuperscript{(44)} సంగతీ అంతే — వాళ్ళు ఉదయం నుంచీ తమ
\VUgras{గాలపు కర్రలతో} \textsuperscript{(45)} నీటి ఒడ్డున వచ్చి కూర్చున్నారు.

%%K t08-c2-09-2 p
\VUgras{చేప} \textsuperscript{(47)} \VUgras{గాలాన్ని} నోట కరుస్తుంది — వాళ్ళు తమ \VUgras{దారం}
\textsuperscript{(46)} నీళ్ళలో వేయగానే; \VUgras{ఎర} మళ్ళీ పెట్టేందుకు వాళ్ళకు తీరిక దొరికిందో
లేదో కొత్త బలి దారం చివర గిలగిలలాడుతుంది.

%%K t08-c2-09-3 p
అయినా \VUgras{వల వేసేవాడు} \textsuperscript{(48)} ఎక్కువ చేపలు పడతాడు — అతను తన
\VUgras{పడవలోంచి} \textsuperscript{(49)} \VUgras{నది} \textsuperscript{(50)} మధ్యలోకి
\VUgras{వల} విసురుతాడు.

%%K t08-c2-10-1 p
10. --- శరదృతువు మంచి విహారాల కాలం : కాలినడకన, \VUgras{గుర్రంపై} \textsuperscript{(52)},
\VUgras{సైకిల్‌పై} \textsuperscript{(53)}, \VUgras{మోటారులో} \textsuperscript{(54)}.
\VUgras{దారులు} \textsuperscript{(51)} శుభ్రంగా ఉన్నాయి, చివరి మంచి రోజులను జనం వాడుకుంటారు;
ఎందుకంటే ఇదిగో \VUgras{స్వాలో పక్షులు} \textsuperscript{(56)} ఇప్పటికే కప్పులపై
గుమిగూడుతున్నాయి, మరింత వెచ్చని దేశాలకు రెక్కలు విప్పి ఎగిరిపోయేందుకు. ముసలి
\VUgras{అక్రోట్ చెట్టు} \textsuperscript{(57)} ఆకులు పసుపెక్కుతున్నాయి, \VUgras{అక్రోట్లు}
రాలుతున్నాయి, \VUgras{పూలగుత్తిలా} \textsuperscript{(58)} గుంపులుగా నిలబడిన ఎత్తైన
\VUgras{చెట్లు} — \VUgras{కొండపైవి} \textsuperscript{(59)} — త్వరలో ఆకులు రాల్చుకుంటాయి.

%%K t08-c2-11-1 p
11. --- \VUgras{సూర్యుడు} \textsuperscript{(60)} ఆలస్యంగా ఉదయిస్తాడు, పగలు తగ్గిపోతుంది,
\VUgras{ఉదయాన} బాగా కోసే చలి తగలడం మొదలవుతుంది; ఇదిగో \VUgras{వుడ్‌కాక్}
\textsuperscript{(61)} గుంపులు ఇప్పటికే మన ఆతిథ్యం లేని వాతావరణాన్ని వదిలి వెళ్తున్నాయి.

\VUsaut{6.0mm}
\VUfilet{22mm}
